\pagestyle{empty}
\tinyskip
\vfill
\begin{flushleft}
\small{
Copyright \copyright{} 2017 Wessmans Musikförlag AB \\
Notskrivning och layout: Erik Ronström \\
Tryck: Wessmans Musikförlag AB, Visby

\vspace{3mm}

Noterna satta med Lilypond 2.19.54 \\
Texten satt med \LaTeX{} (LuaTeX 0.95.0)

\vspace{3mm}

ISMN: 979-0-66163-536-9 \\
ISBN: 978-91-8771-060-5 \\
Beställningsnummer: 201749A
}
\end{flushleft}

\newgeometry{top=33mm,left=36mm,right=36mm,bottom=33mm,twoside,bindingoffset=8mm,headsep=1.8cm}
\fancyhfoffset[E,O]{0pt}
\addtolength{\skip\footins}{\baselineskip}
{
\setlength{\parindent}{1.5em}


%%%%%%%%%%%%%%%%%%%%%%%%%%%%%%%%%%%%%%%%%%%%%%%%%%%%%%%%%%%%%%%%%%%%%%%%%%%%%%%% 

\section*{\centering \LARGE Förord.}\vspace{1cm}
% \markboth{\MakeUppercase{Förord}}{\MakeUppercase{Förord}}

\hyphenation{Gotlands-toner}
\hyphenation{hand-på-lägg-ning}

Gotlandstoner, August Fredins samling av 727 visor och låtar från Gotland,
har en lång historia. Uppteckningsarbetet påbörjades redan 1875 och pågick
i nästan 25 år. Utgivningen kom att ta ännu längre tid och inte förrän 1933  % längre tid. Inte förrän
fanns verket för första gången utgivet i sin helhet.

Min egen relation till Gotlandstoner inleddes 1999. Jag hade då spelat gitarr i
några år, med svensk folkmusik som \guillemotright{}hemmaplan\guillemotright{}, men jag hade också nyligen börjat
vända mina ögon och öron mot Irland. Vi hade ett par skivor hemma med irländsk
folkmusik och dem kunde jag till slut fram- och baklänges. I jakt på mer material letade jag
mig ut på nätet och upptäckte ganska snart att där fanns hemsidor med omfattande samlingar
av irländska låtar, mest jigs och reels. Men låtarna var inte skrivna med noter, utan
med bokstäver och siffror! \textit{abc-formatet}, som det kallas, är praktiskt, eftersom man
enkelt kan skriva musik med tangentbordet, och t.\,ex.\@ e-posta låtar till varandra.

Jag överväldigades av mängden material, men datanörd som jag var blev jag också
fascinerad av abc-formatet i sig. Jag fick lust att prova att själv sätta
ihop en abc-låtsamling. Ganska snart kom jag fram till mitt drömprojekt:
Gotlandstoner, sedan länge ur tryck och dessutom med musik härifrån Gotland!
Att digitalisera den vore ett perfekt projekt! Jag satte igång direkt och skrev väl
in ett tjugotal låtar innan jag tröttnade, och boken hamnade i hyllan igen.

Något år senare gick jag på Malungs folkhögskola och då tog jag tag i Gotlandstoner igen.
Den här gången var jag dock inte ensam: nu delade jag projektet med Clara Andermo.
Under ett par års tid gjorde vi då och då ett ryck och skrev in en bunt låtar. 2004
fanns till slut alla låtar digitalt. Vi publicerade allt material på en hemsida, gotlandstoner.se,
samtidigt som Kulturföreningen Roxy gav ut en faksimil-utgåva av Gotlandstoner.

Vi diskuterade också möjligheten att göra en mer \guillemotright{}spelbar\guillemotright{}
utgåva av boken. Från abc-källan borde det ju vara möjligt att mer eller mindre
automatiskt tillverka noter för tryck! Jag gjorde några provsidor, men det visade sig
krävas mycket mer \guillemotright{}handpåläggning\guillemotright{} än jag trott och idén rann ut i sanden.

Nu har det gått ytterligare 13 år. Under den tiden är det många som frågat om inte
den där ringpärmsutgåvan är på gång snart. Nu äntligen får jag nöjet att svara ja på den frågan!

\bigskip

{
\setlength{\parindent}{0em}
Denna utgåva utgör alltså det hittills sista kapitlet i
Gotlandstoners långa historia. Jag vill tacka alla er som kommit med hejarop längs vägen
och som fortsatt fråga --- utan er uppmuntran hade det aldrig blivit färdigt! Allra mest
tack skyldig är jag min pappa Owe, som ständigt stöttat mig, såväl i detta projekt som i allt annat.
}

Jag vill också rikta ett särskilt tack till Jan Lundqvist och Fredrik Lönngren för ett omfattande
korrekturarbete, till Örjan Klintberg på Wessmans Musikförlag samt till Kungl.\@ Gustaf Adolfs Akademien
för ekonomiskt bidrag till projektet.

\bigskip

\begin{flushright}
\textit{Erik Ronström} \\
\textit{Mästerby, juni 2017}
\end{flushright}


\clearpage

%%%%%%%%%%%%%%%%%%%%%%%%%%%%%%%%%%%%%%%%%%%%%%%%%%%%%%%%%%%%%%%%%%%%%%%%%%%%%%%% 

\section*{\centering \LARGE Om denna utgåva.}\vspace{8mm}
% \markboth{\MakeUppercase{Förord}}{\MakeUppercase{Förord}}


% Målet med denna utgåva är att den i första hand ska vara \textit{användbar}
% att spela efter.

\makeatletter
\renewcommand\paragraph{\@startsection{paragraph}{4}{\z@}%
                                    {3.0ex \@plus1ex \@minus.2ex}%
                                    {-1em}%
                                    {\normalfont\normalsize\bfseries}}

\makeatother

Utgångspunkten för denna utgåva är att återge originalutgåvan i en ny läsligare
form, men med oförändrat innehåll. I några avseenden skiljer sig dock utgåvorna åt.

\paragraph{Texter.} Originalutgåvan innehåller en över 50 sidor lång inledning,
som ger en kulturhistorisk bakgrund och inramning till det musikaliska materialet.
Där finns inte minst en \guillemotright{}Förteckning över utövare av gotländsk
folkmusik\guillemotright{}.
Av utrymmesskäl har vi dock valt att i denna utgåva bara ta med sådana texter
som direkt relaterar till låtarna och hur de spelats. Dessa avsnitt
har placerats i anslutning till respektive låttyp.

För den som vill läsa mer av Gotlandstoners texter hänvisas till originalutgåvan,
och till hemsidan \href{http://www.gotlandstoner.se/}{www.gotlandstoner.se}, där textmaterialet återfinns i sin helhet.
På hemsidan finns också det förord som Owe Ronström och
Märta Ramsten skrev till Gotlandstoners faksimilutgåva 2004, och som
sätter Gotlandstoner och August Fredin i ett historiskt sammanhang.

\paragraph{Fonetisk skrift.} Tecknet \raisebox{-0.3\height}{\includegraphics{include/g10-crop.pdf}} används
på vissa ställen i originalutgåvan för att beteckna ng-ljudet, det är i denna utgåva ersatt med tecknet \ng{}.

\paragraph{Korrigeringar.} De rättelser som finns längst bak i originalutgåvan har
införts i texten och noterna i denna utgåva. Därutöver har minimala redigeringar
gjorts på ett antal ställen (mindre stav- eller tryckfel, felaktiga sidhänvisningar etc).
\guillemotright{}Musikaliska fel\guillemotright{} har rättats utan särskild notis
om de varit \guillemotright{}triviala\guillemotright{}, t.\,ex.\@ när en takt
varit av fel längd på grund av en saknad balk, och det är uppenbart hur det borde ha varit.

Det finns också många ställen i originalutgåvan där man kan \textit{misstänka}
tryckfel i noterna. I dessa fall har generellt inte gjorts någon ändring, och
om någon redigering gjorts är denna tydligt utmärkt.

\paragraph{Balkning.} Balkningen är i denna utgåva omgjord efter den nutida principen
att balkarna i första hand visar slagindelningen.

\begingroup
\widowpenalty 10000
\clubpenalty10000
\paragraph{Taktarter.} Ett fåtal låtar är i originalutgåvan noterade i \sfrac{3}{8},
de har ändrats till \sfrac{3}{4}.
\endgroup

\paragraph{Scordatura.} En del av låtarna är upptecknade efter omstämd fiol, och några av
dessa är i originalutgåvan skrivna i s.\,k.\@ \textit{scordatura}-notation, där noterna
motsvarar fingersättningen på fiolen snarare än de toner som klingar. I denna
utgåva är alla låtar skrivna klingande.

\paragraph{Reprisering.} På några ställen förekommer i originalutgåvan följande mönster:

\vspace{3mm}
\includegraphics{include/snippets/repriser-fel-crop.pdf}
\vspace{3mm}

{
\setlength{\parindent}{0em}
Det är svårt att göra en entydig tolkning: vart går man när man spelat genom det
andra \guillemotright{}huset\guillemotright{} (2:an)? Jag har valt att tolka det som ett enkelt tryckfel, och
i stället skrivit:
}

\vspace{3mm}
\includegraphics{include/snippets/repriser-ratt-crop.pdf}

}
\restoregeometry
\fancyhfoffset[E,O]{0pt}
